\chapter{DHCP, NAT e IPv6}

% ══════════════════════════════════════════════════════════════
\section{L'assegnazione degli indirizzi IP}
% ══════════════════════════════════════════════════════════════

Dentro un blocco di indirizzi, gli indirizzi IP vanno assegnati alle singole interfacce. Un amministratore di sistema
può farlo in due modi:

\begin{itemize}
    \item \textbf{manualmente}: assegnando a mano un indirizzo a ciascun host (indirizzo \emph{statico});
    \item \textbf{automaticamente}: è la rete ad assegnare un indirizzo libero agli host che arrivano.
\end{itemize}

Il secondo approccio, il più comune, si realizza con il \textbf{Dynamic Host Configuration Protocol} (DHCP).

\dfn{DHCP}{
    Il \textbf{DHCP} assegna automaticamente a un host che si collega alla rete le informazioni necessarie per
    usarla: oltre all'\textbf{indirizzo IP}, la \textbf{subnet mask}, l'indirizzo del \textbf{gateway predefinito}
    (per uscire dalla rete) e quello del \textbf{server DNS locale}. È \emph{plug-and-play}, ed è quello che si usa
    nelle nostre case.
}

% ══════════════════════════════════════════════════════════════
\section{Il funzionamento del DHCP}
% ══════════════════════════════════════════════════════════════

Il DHCP è un protocollo \textbf{client-server}: un host appena arrivato (client) contatta il server DHCP per ricevere
la configurazione di rete. Il servizio può essere offerto da un computer o dal router stesso. Formalmente, il DHCP è un
protocollo di \textbf{livello applicativo} (usa UDP).

Riprendiamo l'esempio delle tre subnet collegate da un router, e supponiamo che la subnet di destra
(\texttt{223.1.2.0/24}) abbia un server DHCP: quando un nuovo host entra nella rete, ``contratta'' il proprio indirizzo
con il server. L'aggiunta di un host avviene in quattro passi (ricordati spesso con l'acronimo \textbf{DORA}):

\begin{figure}[H]
\centering
\begin{minipage}{0.34\linewidth}\centering
  \includegraphics[width=0.9\linewidth]{cap15/dhcp_messaggi.png}
\end{minipage}\hfill
\begin{minipage}{0.62\linewidth}
\begin{enumerate}
  \item \textbf{Discover} --- ricerca del server: il nuovo host invia in \textbf{broadcast} un messaggio \emph{DHCP discover}
        (UDP, porta \textbf{67}) per trovare un server DHCP.
        \begin{itemize}
          \item IP destinazione: \texttt{255.255.255.255} (broadcast);
          \item IP sorgente: \texttt{0.0.0.0} (``questo host'', che non ha ancora un indirizzo);
          \item un identificativo di transazione casuale.
        \end{itemize}
  \item \textbf{Offer} --- offerta: poiché possono esserci più server, ogni server che riceve il discover risponde con un
        messaggio in broadcast (porta \textbf{68}) che offre una possibile configurazione (campo \texttt{yiaddr}, \emph{Your
        Internet ADDRess}). IP destinazione \texttt{255.255.255.255}, IP sorgente quello del server (\texttt{223.1.2.5}).
  \item \textbf{Request} --- richiesta: il client sceglie un'offerta tra quelle ricevute e la accetta ``facendo l'eco''
        del messaggio di offerta.
  \item \textbf{ACK} --- conferma: il server conferma la transazione con un messaggio finale.
\end{enumerate}
\end{minipage}
\caption{I quattro messaggi del DHCP (da Kurose).}
\end{figure}

\warning{Perché tutto in broadcast}{
    Nel passo 1 il client \textbf{non ha ancora un indirizzo} e \textbf{non sa dove sia il server}: l'unico modo di
    raggiungerlo è il broadcast. Anche l'offerta viaggia in broadcast, perché il client non può ancora ricevere
    pacchetti indirizzati a un IP che non ha. L'indirizzo assegnato ha una \textbf{durata} (\emph{lease time}): scaduta,
    va rinnovata, altrimenti può essere riassegnata a un altro host.
}

\subsection{Il DHCP relay agent}

\begin{itemize}
    \item Nell'esempio tutto funziona perché il server DHCP e l'host sono nella \textbf{stessa subnet}: i broadcast non
          attraversano i router.
    \item Se il server esiste ma sta in un'altra subnet, serve un \textbf{DHCP relay agent} che inoltri i messaggi DHCP (può
          farlo un router).
    \item Configurando il router come relay agent, anche gli host delle subnet \texttt{223.1.1.0/24} e \texttt{223.1.3.0/24}
          vengono serviti dal server DHCP della subnet \texttt{223.1.2.0/24}.
\end{itemize}

\subsection{Verificare la configurazione}

Su Linux il comando \texttt{ip} mostra la configurazione ricevuta dal DHCP:

\begin{lstlisting}[language=bash]
$ ip addr       # indirizzo IP e maschera delle interfacce
$ ip route      # gateway predefinito (riga "default via ...")
\end{lstlisting}

% ══════════════════════════════════════════════════════════════
\section{Il problema della visibilità}
% ══════════════════════════════════════════════════════════════

Su Internet ci sono miliardi di dispositivi che dovrebbero avere un indirizzo IP univoco per essere raggiungibili da
chiunque. Ma è davvero necessario che \textbf{tutti} i dispositivi siano visibili a tutti? È una buona idea che la mia
stampante sia raggiungibile (e usabile) da chiunque su Internet?

Avere tutti i dispositivi raggiungibili dall'esterno delle reti locali pone problemi evidenti:

\begin{itemize}
    \item gli indirizzi IP sono \textbf{finiti} (e gli IPv4 praticamente esauriti);
    \item se i dispositivi sono collegati localmente, spesso è irragionevole o indesiderabile esporli tutti su Internet;
    \item gli amministratori locali dovrebbero conoscere la struttura dei blocchi di indirizzi soprastanti per assegnare
          indirizzi liberi, cosa che spesso non dipende da loro (per esempio nelle reti di casa).
\end{itemize}

% ══════════════════════════════════════════════════════════════
\section{Il NAT}
% ══════════════════════════════════════════════════════════════

\dfn{Network Address Translation (NAT)}{
    Il servizio di \textbf{NAT} permette di \textbf{rimappare} gli indirizzi IP dei pacchetti di una rete in indirizzi
    diversi (\emph{network masquerading}). Si usa una \textbf{tabella di traduzione} che associa molte coppie
    indirizzo/porta locali (della LAN) a una coppia indirizzo/porta globale (della WAN). Il servizio può essere offerto
    dai router.
}

La rete locale mascherata si chiama \textbf{rete privata} (o \emph{realm} con indirizzi privati): gli indirizzi privati
sono validi \textbf{solo all'interno} della rete privata.

\begin{figure}[H]
\centering
\begin{tikzpicture}[every node/.style={font=\footnotesize\sffamily}]
  \node (pc) at (0,1.4) {\icon[0.9cm]{pc}}; \node[left=0pt of pc, font=\scriptsize\ttfamily] {192.168.0.23};
  \node (lp) at (0,0) {\icon[1cm]{laptop}}; \node[left=0pt of lp, font=\scriptsize\ttfamily] {192.168.0.11};
  \node (ph) at (0,-1.4) {\icon[0.6cm]{smartphone}}; \node[left=0pt of ph, font=\scriptsize\ttfamily] {192.168.0.7};
  \node (r) at (4,0) {\icon[1.3cm]{router}}; \node[below=0pt of r, align=center] {router NAT};
  \foreach \n in {pc,lp,ph} \draw[wired] (\n) -- (r);
  \node[netcloud, minimum width=3cm] (i) at (9,0) {Internet};
  \draw[wired, line width=2pt] (r) -- node[above, font=\scriptsize\ttfamily] {203.0.113.221} (i);
  \draw[decorate, decoration={brace, amplitude=6pt, mirror}] (-1.8,-2.1) -- (4.6,-2.1) node[midway, below=7pt, text=netgreen!50!black] {rete privata: indirizzi 192.168.0.x};
  \draw[decorate, decoration={brace, amplitude=6pt, mirror}] (4.8,-2.1) -- (10.6,-2.1) node[midway, below=7pt, text=netblue] {per Internet c'è un solo indirizzo pubblico};
\end{tikzpicture}
\caption{Con il NAT tutti i dispositivi di casa escono verso Internet con lo stesso indirizzo pubblico del router.}
\end{figure}

\subsection{Un esempio di NAT}

\begin{figure}[H]
    \centering
    \includegraphics[width=0.7\linewidth]{cap15/nat_esempio.png}
    \caption{Un router NAT tra una rete privata 10.0.0.0/24 e Internet (da Kurose).}
\end{figure}

Le quattro interfacce della rete locale hanno indirizzi della subnet \texttt{10.0.0.0/24}; lato Internet il router ha un solo
indirizzo, \texttt{138.76.29.7}. Il router si comporta da ``passamessaggi'' che sovrascrive gli indirizzi secondo la tabella:

\begin{enumerate}
    \item L'host \texttt{10.0.0.1} invia un datagramma verso il server Web \texttt{128.119.40.186}, porta 80, usando la porta sorgente 3345.
    \item Il router NAT riceve il datagramma, sceglie una porta libera (\textbf{5001}), \textbf{sostituisce} indirizzo e porta sorgente
          con \texttt{138.76.29.7, 5001} e aggiunge la riga alla tabella di traduzione.
    \item Il server risponde a \texttt{138.76.29.7, 5001}: per lui, la richiesta è arrivata dal router.
    \item Il router riceve la risposta, cerca \texttt{138.76.29.7, 5001} nella tabella e \textbf{riscrive} destinazione e porta con
          \texttt{10.0.0.1, 3345}, poi la inoltra nella LAN.
\end{enumerate}

\begin{table}[H]
\centering
\begin{tabular}{ll}
\toprule
\textbf{Lato WAN} & \textbf{Lato LAN} \\
\midrule
138.76.29.7, 5001 & 10.0.0.1, 3345 \\
138.76.29.7, 5002 & 10.0.0.2, 3345 \\
\dots & \dots \\
\bottomrule
\end{tabular}
\caption{La tabella di traduzione NAT: la seconda riga mostra perché serve una porta nuova.}
\end{table}

\warning{Perché il NAT cambia anche la porta}{
    Gli host non sanno nulla del traffico degli altri host: due host della LAN potrebbero scegliere \textbf{la stessa
    porta} sorgente (3345) nello stesso momento. Il NAT non può quindi riusare la porta iniziale, e ne assegna una nuova.
    Poiché la porta è di 16 bit, un NAT può gestire \textbf{oltre 60\,000 connessioni} contemporanee con un solo indirizzo IP pubblico.
}

\subsection{Il limite del NAT}

\info{NAT traversal}{
    Nell'esempio la comunicazione è stata \textbf{iniziata dall'host interno} \texttt{10.0.0.1}: la riga della tabella nasce
    proprio da quel primo pacchetto in uscita. Se un pacchetto arriva da Internet \textbf{senza comunicazioni precedenti}
    (qualcuno vuole contattare un server che sta dietro il NAT), il router non ha informazioni sufficienti per sapere a quale
    host interno consegnarlo. Servono allora tecniche aggiuntive, dette di \textbf{NAT traversal} (per esempio l'inoltro
    manuale delle porte sul router, o protocolli per scoprire e aprire le mappature). Il NAT è anche criticato perché un router,
    dispositivo di livello 3, modifica i numeri di porta, che appartengono al livello 4.
}

% ══════════════════════════════════════════════════════════════
\section{Gli indirizzi IP riservati}
% ══════════════════════════════════════════════════════════════

Perché gli amministratori locali possano organizzare una rete privata senza interferenze, per convenzione alcuni blocchi di
indirizzi sono \textbf{riservati}: ISP e ICANN non possono assegnarli a nessuna organizzazione.

\begin{table}[H]
\centering
\begin{tabular}{llp{7.5cm}}
\toprule
\textbf{Blocco} & \textbf{Intervallo} & \textbf{Uso} \\
\midrule
10.0.0.0/8 & 10.0.0.0 -- 10.255.255.255 & reti private \\
172.16.0.0/12 & 172.16.0.0 -- 172.31.255.255 & reti private \\
192.168.0.0/16 & 192.168.0.0 -- 192.168.255.255 & reti private (tipiche delle case) \\
127.0.0.0/8 & 127.0.0.0 -- 127.255.255.255 & ``questa macchina'' (\emph{loopback}), di solito 127.0.0.1 \\
0.0.0.0 & --- & rete corrente; per estensione, un indirizzo con la parte di host a 0 indica la rete \\
255.255.255.255 & --- & broadcast; per estensione, un indirizzo con la parte di host tutta a 1 è il broadcast della subnet \\
\bottomrule
\end{tabular}
\caption{Indirizzi IPv4 riservati.}
\end{table}

La riserva è necessaria perché gli indirizzi privati potrebbero coincidere con indirizzi pubblici, rendendo ambiguo
l'instradamento dentro la rete. Con blocchi riservati, un router sa che \texttt{192.168.x.x} non è mai un indirizzo di Internet.

% ══════════════════════════════════════════════════════════════
\section{ping e nmap per la scoperta degli host}
% ══════════════════════════════════════════════════════════════

Per verificare se un host è raggiungibile si usa \textbf{ping} (identico su Linux e Windows), basato su \textbf{ICMP}
(\emph{Internet Control Message Protocol}): invia all'host un pacchetto \emph{echo request} a cui l'host risponde
automaticamente con un \emph{echo reply}, e misura il tempo trascorso, cioè l'\textbf{RTT}. Lo vedremo in dettaglio nel
prossimo capitolo.

Oltre alla scansione delle porte, \textbf{nmap} può mappare i dispositivi di una rete, usando ping (ICMP) per scoprire gli host:

\begin{lstlisting}[language=bash]
$ ping [indirizzo]                               # l'host risponde?
$ sudo nmap -sn [indirizzo gateway]/[bit di subnet]  # quali host sono accesi nella subnet?
\end{lstlisting}

% ══════════════════════════════════════════════════════════════
\section{Esempio completo: configurare una LAN}
% ══════════════════════════════════════════════════════════════

Vogliamo creare una nuova rete locale (LAN) dentro una rete più grande preesistente (WAN). L'amministratore della WAN
ci fornisce:

\begin{itemize}
    \item l'indirizzo della WAN: \texttt{150.100.0.0/16};
    \item l'indirizzo del gateway che porta all'esterno (WAN/Internet): \texttt{150.100.50.1};
    \item un indirizzo libero per la nostra rete: \texttt{150.100.50.10}.
\end{itemize}

Useremo un router con NAT e decidiamo che la nuova LAN avrà indirizzo \texttt{192.168.1.0/24} (riservato alle reti private).

\begin{figure}[H]
\centering
\begin{tikzpicture}[every node/.style={font=\footnotesize\sffamily}]
  \node[netcloud, minimum width=2.8cm] (inet) at (0,3.2) {Internet};
  \node[draw=netblue, thick, rounded corners, fill=cloudbg, minimum width=7cm, minimum height=2.2cm] (wan) at (0,0.8) {};
  \node[anchor=north east, text=netblue] at (wan.north east) {WAN \texttt{150.100.0.0/16}};
  \node (rw) at (-1.8,1.0) {\icon[1.1cm]{router}}; \node[below=-1pt of rw, align=center, font=\scriptsize\sffamily] {router WAN\\\texttt{150.100.50.1}};
  \draw[wired, line width=2pt] (inet) -- (rw);
  \node[draw=netgreen!70!black, thick, rounded corners, fill=netgreen!8, minimum width=7cm, minimum height=2.6cm] (lan) at (0,-2.4) {};
  \node[anchor=south east, text=netgreen!50!black] at (lan.south east) {LAN \texttt{192.168.1.0/24}};
  \node (rl) at (-1.8,-1.6) {\icon[1.1cm]{router}};
  \node[anchor=west, font=\scriptsize\ttfamily, text=netblue] at (-1.2,-0.9) {WAN: 150.100.50.10};
  \node[anchor=west, font=\scriptsize\ttfamily, text=netgreen!50!black] at (-1.2,-1.8) {LAN: 192.168.1.1};
  \draw[wired] (rw) -- (rl);
  \node (h1) at (-2.3,-3.2) {\icon[0.8cm]{laptop}}; \node[right=0pt of h1, font=\scriptsize\ttfamily, align=left] {192.168.1.100\\(DHCP)};
  \node (h2) at (1.2,-3.2) {\icon[0.8cm]{pc}}; \node[right=0pt of h2, font=\scriptsize\ttfamily, align=left] {192.168.1.10\\(manuale)};
  \draw[wired] (rl) -- (h1); \draw[wired] (rl) -- (h2);
\end{tikzpicture}
\caption{La LAN da configurare: il router locale ha un'interfaccia nella WAN e una nella LAN.}
\end{figure}

\subsection{1. Il lato WAN del router}

Il router locale ha \textbf{due interfacce}, una verso la WAN e una verso la LAN. Sul lato WAN vanno impostati:

\begin{itemize}
    \item l'indirizzo del router nella WAN: \texttt{150.100.50.10};
    \item la subnet mask della WAN: \texttt{255.255.0.0};
    \item il gateway: \texttt{150.100.50.1};
    \item uno o più server DNS.
\end{itemize}

Se la WAN ha un servizio DHCP, si può anche scegliere un indirizzo dinamico.

\begin{figure}[H]
    \centering
    \includegraphics[width=0.6\linewidth]{cap15/config_wan.png}
    \caption{La pagina di configurazione ``Internet'' (lato WAN) di un router domestico: indirizzo statico, maschera, gateway e DNS (valori di un esempio reale, diversi da quelli del testo).}
\end{figure}

\subsection{2. Il lato LAN del router}

\begin{itemize}
    \item All'interfaccia LAN del router si dà un indirizzo della nuova rete: poiché la LAN è \texttt{192.168.1.0/24}, al router si
          assegna \texttt{192.168.1.1/24} (è consuetudine dare al router l'indirizzo \texttt{.1}).
    \item Se il router offre anche il DHCP, si specificano:
    \begin{itemize}
      \item il \textbf{pool di indirizzi} per gli host DHCP, per esempio da \texttt{192.168.1.100} a \texttt{192.168.1.250};
      \item il \textbf{lease time}: la durata dell'assegnazione, dopo la quale l'indirizzo può essere riusato;
      \item gateway e DNS da comunicare agli host.
    \end{itemize}
\end{itemize}

\begin{figure}[H]
    \centering
    \includegraphics[width=0.6\linewidth]{cap15/config_lan.png}
    \caption{La configurazione del server DHCP del router: pool di indirizzi, lease time, gateway e DNS.}
\end{figure}

\subsection{3. Gli host}

Su un nuovo host (per esempio con Linux Mint) si può scegliere come configurare la connessione:

\begin{figure}[H]
\centering
\begin{minipage}{0.44\linewidth}\centering
  \includegraphics[width=0.85\linewidth]{cap15/host_dhcp.png}\\{\small (a) automatica (DHCP)}
\end{minipage}\hfill
\begin{minipage}{0.44\linewidth}\centering
  \includegraphics[width=0.85\linewidth]{cap15/host_manuale.png}\\{\small (b) manuale}
\end{minipage}
\caption{Configurazione della connessione su un host.}
\end{figure}

\begin{itemize}
    \item \textbf{Automatica (DHCP)}: l'host riceve dal DHCP un indirizzo tra \texttt{192.168.1.100} e \texttt{192.168.1.250}.
    \item \textbf{Manuale}: bisogna conoscere la configurazione della rete e gli indirizzi disponibili, e indicare
          l'\textbf{indirizzo dell'host} (statico, \textbf{fuori dal pool} del DHCP, per esempio \texttt{192.168.1.10}), la
          \textbf{subnet mask}, il \textbf{gateway} e uno o più \textbf{DNS}.
\end{itemize}

Gli host collegati inviano i loro pacchetti al router locale, senza sapere nulla della WAN esterna; il router locale è
configurato per inoltrarli al router della WAN, e da lì eventualmente a Internet.

\warning{Indirizzi manuali fuori dal pool}{
    Un indirizzo assegnato a mano deve stare \textbf{fuori} dall'intervallo del DHCP: altrimenti prima o poi il DHCP lo
    assegnerà anche a un altro host, e due interfacce avranno lo stesso indirizzo (conflitto).
}

% ══════════════════════════════════════════════════════════════
\section{IPv6}
% ══════════════════════════════════════════════════════════════

All'inizio degli anni '90 l'IETF (organizzazione nata nel 1986) avviò lo sviluppo di un successore di IPv4, per
affrontare la \textbf{carenza di indirizzi} a 32 bit. IPv6 fu progettato per avere molti più indirizzi, ma fu anche
l'occasione per aggiornare altri aspetti di IPv4 sulla base dell'esperienza accumulata.

\info{E IPv5?}{
    Non si sa esattamente quando gli IPv4 si esauriranno del tutto (si era ipotizzato il 2008, poi il 2018, e qualche
    indirizzo è ancora disponibile). \textbf{IPv5} fu proposto nel 1979 e si basava sullo sperimentale \emph{Internet Stream
    Protocol} (ST); usava ancora indirizzi a 32 bit, e fu abbandonato soprattutto per la stessa carenza di indirizzi.
}

\subsection{I vantaggi di IPv6}

\begin{itemize}
    \item \textbf{Indirizzi molto più grandi}: da 32 a \textbf{128 bit}, cioè $2^{128} \approx 3{,}4\cdot 10^{38}$ indirizzi. Il mondo
          non resterà senza indirizzi: si potrebbe dare un indirizzo a ogni granello di sabbia del pianeta.
    \item \textbf{Header di lunghezza fissa}, 40 byte: alcuni campi di IPv4 sono stati eliminati o resi facoltativi, e i router
          elaborano i pacchetti più velocemente.
    \item \textbf{Etichettatura dei flussi} (\emph{flow labeling}): i pacchetti di certe applicazioni (per esempio lo streaming
          audio/video) possono essere raggruppati in un flusso con un identificativo. Il ruolo preciso dei flussi è ancora discusso.
\end{itemize}

\info{Come si scrive un indirizzo IPv6}{
    Un indirizzo IPv6 si scrive come 8 gruppi di 4 cifre esadecimali separati da due punti, per esempio
    \texttt{2001:0db8:0000:0000:0000:ff00:0042:8329}. Gli zeri iniziali di ogni gruppo si possono omettere, e una sequenza di
    gruppi tutti a zero si può sostituire (una sola volta) con \texttt{::}: l'indirizzo diventa \texttt{2001:db8::ff00:42:8329}.
    L'indirizzo di loopback è \texttt{::1}.
}

\subsection{Il datagramma IPv6}

\begin{figure}[H]
    \centering
    \includegraphics[width=0.55\linewidth]{cap15/ipv6_datagramma.png}
    \caption{Il datagramma IPv6: header fisso di 40 byte.}
\end{figure}

\begin{itemize}
    \item \textbf{Version} (4 bit): la versione, cioè 6.
    \item \textbf{Traffic class} (8 bit): come il TOS di IPv4, serve a dare priorità ai datagrammi (per esempio voce su IP rispetto a SMTP).
    \item \textbf{Flow label} (20 bit): identifica un flusso di datagrammi.
    \item \textbf{Payload length} (16 bit): numero di byte del payload (esclusi i 40 byte di header).
    \item \textbf{Next header} (8 bit): il protocollo superiore a cui consegnare i dati (TCP, UDP, \dots), come il campo protocollo di IPv4.
    \item \textbf{Hop limit} (8 bit): il time-to-live, decrementato a ogni inoltro.
    \item \textbf{Indirizzi sorgente e destinazione} (128 + 128 bit).
    \item \textbf{Data}: il payload.
\end{itemize}

\subsection{Cosa manca rispetto a IPv4}

\begin{itemize}
    \item \textbf{Campi della frammentazione}: IPv6 non permette frammentazione e riassemblaggio nei router intermedi, ma solo negli
          host. Se un router riceve un datagramma IPv6 troppo grande da inoltrare, lo \textbf{scarta} e rimanda al mittente un
          messaggio di errore ``\emph{Packet Too Big}''; il mittente rinvia i dati con datagrammi più piccoli.
    \item \textbf{Checksum dell'header}: richiede tempo in ogni router ed è ridondante, perché i protocolli di collegamento di
          solito controllano già l'intero pacchetto (e TCP/UDP controllano i propri dati).
    \item \textbf{Opzioni}: non fanno più parte dell'header standard; alcune si possono indicare nel campo \emph{next header}
          (accanto agli identificativi di TCP/UDP), come header di estensione.
\end{itemize}

% ══════════════════════════════════════════════════════════════
\section{Dalla IPv4 alla IPv6}
% ══════════════════════════════════════════════════════════════

C'è un grosso problema: IPv6 \textbf{non è retrocompatibile}. I sistemi che capiscono solo IPv4 non sanno gestire i datagrammi IPv6.
Come può Internet adottare IPv6?

\begin{itemize}
    \item \textbf{Flag day}: a una certa data e ora tutte le macchine di Internet vengono spente e aggiornate da IPv4 a IPv6. Un
          approccio simile fu usato quando fu introdotto TCP, e fu un disastro anche nella piccola Internet di allora: con miliardi
          di dispositivi è impensabile.
    \item \textbf{Tunneling}: router specifici (i \emph{tunnel}) incapsulano i datagrammi IPv6 dentro datagrammi IPv4 (e viceversa)
          in modo quasi trasparente, per attraversare le parti di rete che parlano solo IPv4. È l'approccio più realistico, già
          usato nella pratica.
\end{itemize}

\begin{figure}[H]
\centering
\begin{tikzpicture}[every node/.style={font=\scriptsize\sffamily}]
  \foreach \x/\n/\c in {0/A/lnet, 2.6/B/lnet, 5.2/C/lphy, 7.8/D/lphy, 10.4/E/lnet, 13/F/lnet} {
    \node (\n) at (\x,0) {\icon[0.9cm]{router}}; \node[above=-2pt of \n] {\n};
  }
  \node[below=0pt of A] {IPv6}; \node[below=0pt of B] {IPv6/v4}; \node[below=0pt of C] {IPv4}; \node[below=0pt of D] {IPv4}; \node[below=0pt of E] {IPv6/v4}; \node[below=0pt of F] {IPv6};
  \foreach \a/\b in {A/B, B/C, C/D, D/E, E/F} \draw[wired] (\a) -- (\b);
  \draw[very thick, netorange, rounded corners, dashed] (2.4,-1.2) rectangle (10.6,0.9);
  \node[text=netorange, font=\footnotesize\bfseries\sffamily] at (6.5,1.1) {tunnel};
  \node[field, fill=lnet!60, minimum width=1.6cm] at (1.3,-1.8) {datagramma IPv6};
  \node[field, fill=lphy!80, minimum width=1.1cm] (h4) at (5.6,-1.8) {header IPv4};
  \node[field, fill=lnet!60, minimum width=1.6cm, right=0pt of h4] {datagramma IPv6};
  \node[field, fill=lnet!60, minimum width=1.6cm] at (11.7,-1.8) {datagramma IPv6};
\end{tikzpicture}
\caption{Tunneling: B incapsula il datagramma IPv6 in un datagramma IPv4 per attraversare C e D; E lo estrae.}
\end{figure}

L'adozione di IPv6, inizialmente lenta, sta accelerando.

% ══════════════════════════════════════════════════════════════
\section{Riepilogo}
% ══════════════════════════════════════════════════════════════

\riepilogo{
\begin{itemize}
    \item Il \textbf{DHCP} assegna automaticamente indirizzo, maschera, gateway e DNS con quattro messaggi: discover, offer,
          request, ACK (broadcast, porte UDP 67/68); oltre le subnet serve un relay agent.
    \item Il \textbf{NAT} fa uscire una rete privata con un solo indirizzo pubblico, riscrivendo indirizzi e porte secondo una
          tabella; le connessioni avviate dall'esterno richiedono NAT traversal.
    \item Blocchi \textbf{riservati}: 10/8, 172.16/12, 192.168/16 (privati), 127/8 (loopback), 0.0.0.0, 255.255.255.255.
    \item \textbf{IPv6}: indirizzi a 128 bit, header fisso di 40 byte, flow label; niente frammentazione nei router, niente checksum,
          niente opzioni nell'header; transizione con \textbf{tunneling}.
\end{itemize}
}
